\documentclass[10pt,a4paper]{amsart}
\usepackage{amsmath,amssymb,amsthm,graphicx}
% For dealing with references we use the comment environment
\usepackage{verbatim}
\newenvironment{reference}{\comment}{\endcomment}
%\newenvironment{reference}{}{}
\newenvironment{slogan}{\comment}{\endcomment}
\newenvironment{history}{\comment}{\endcomment}
% For commutative diagrams we use Xy-pic
\usepackage[all]{xy}

% We use 2cell for 2-commutative diagrams.
\xyoption{2cell}
\UseAllTwocells
\theoremstyle{plain}
\newtheorem{theorem}[subsection]{Theorem}
\newtheorem{proposition}[subsection]{Proposition}
\newtheorem{lemma}[subsection]{Lemma}

\theoremstyle{definition}
\newtheorem{definition}[subsection]{Definition}
\newtheorem{example}[subsection]{Example}
\newtheorem{exercise}[subsection]{Exercise}
\newtheorem{situation}[subsection]{Situation}

\theoremstyle{remark}
\newtheorem{remark}[subsection]{Remark}
\newtheorem{remarks}[subsection]{Remarks}

\numberwithin{equation}{subsection}
% Macros
%
\def\lim{\mathop{\mathrm{lim}}\nolimits}
\def\colim{\mathop{\mathrm{colim}}\nolimits}
\def\Spec{\mathop{\mathrm{Spec}}}
\def\Hom{\mathop{\mathrm{Hom}}\nolimits}
\def\Ext{\mathop{\mathrm{Ext}}\nolimits}
\def\SheafHom{\mathop{\mathcal{H}\!\mathit{om}}\nolimits}
\def\SheafExt{\mathop{\mathcal{E}\!\mathit{xt}}\nolimits}
\def\Sch{\mathit{Sch}}
\def\Mor{\mathop{\mathrm{Mor}}\nolimits}
\def\Ob{\mathop{\mathrm{Ob}}\nolimits}
\def\Sh{\mathop{\mathit{Sh}}\nolimits}
\def\NL{\mathop{N\!L}\nolimits}
\def\CH{\mathop{\mathrm{CH}}\nolimits}
\def\proetale{{pro\text{-}\acute{e}tale}}
\def\etale{{\acute{e}tale}}
\def\QCoh{\mathit{QCoh}}
\def\Ker{\mathop{\mathrm{Ker}}}
\def\Im{\mathop{\mathrm{Im}}}
\def\Coker{\mathop{\mathrm{Coker}}}
\def\Coim{\mathop{\mathrm{Coim}}}

% Boxtimes
%
\DeclareMathSymbol{\boxtimes}{\mathbin}{AMSa}{"02}

%
% Macros for moduli stacks/spaces
%
\def\QCohstack{\mathcal{QC}\!\mathit{oh}}
\def\Cohstack{\mathcal{C}\!\mathit{oh}}
\def\Spacesstack{\mathcal{S}\!\mathit{paces}}
\def\Quotfunctor{\mathrm{Quot}}
\def\Hilbfunctor{\mathrm{Hilb}}
\def\Curvesstack{\mathcal{C}\!\mathit{urves}}
\def\Polarizedstack{\mathcal{P}\!\mathit{olarized}}
\def\Complexesstack{\mathcal{C}\!\mathit{omplexes}}
% \Pic is the operator that assigns to X its picard group, usage \Pic(X)
% \Picardstack_{X/B} denotes the Picard stack of X over B
% \Picardfunctor_{X/B} denotes the Picard functor of X over B
\def\Pic{\mathop{\mathrm{Pic}}\nolimits}
\def\Picardstack{\mathcal{P}\!\mathit{ic}}
\def\Picardfunctor{\mathrm{Pic}}
\def\Deformationcategory{\mathcal{D}\!\mathit{ef}}
\usepackage{cleveref}
\crefname{originallemma}{original lemma}{original lemmas}
\Crefname{originallemma}{Original lemma}{Original lemmas}
\makeatletter
\newlabel{lemma-CLambdahat-pushouts}{{06SB}{source}{}{sourcepointer-1}{}}
\newlabel{lemma-CLambdahat-pushouts@cref}{{[originallemma][0][]06SB}{[1][0][]source}{}{}{}}
\makeatother
\title{Lemma homotopy in formal-defos.tex}
\author{The Stacks Project Authors}
\date{stacks-project:a04446e57ec1fbc252a871afcec7752fb2807b14; GNU Free Documentation License 1.2 (as published in the Stacks Project COPYING)}
\begin{document}
\maketitle
\noindent\textit{Editorial note.} All mathematical statement and proof text is reproduced verbatim from the author's chapter formal-defos.tex of the Stacks Project, fixed at commit a04446e57ec1fbc252a871afcec7752fb2807b14. The pinned preamble is supplied as its own file; the wrapper class, the xr-hyper wiring and the \textbackslash{}externaldocument list are omitted because the assembler supplies the class.\par\medskip
\noindent\textit{Editorial note.} Numbering inside this standalone excerpt is editorial, not the full Stacks book print numbering.\par\medskip
\noindent\textit{Editorial note.} The candidate is from the corrected band-I follow-on screen (line C) of the extractor's stacks-candidates line: its author label is absent from the live master in every form searched and it belongs to none of the earlier judged pools. The H2 judgement and its reason were made in this module by reading the author's own statement and complete proof in the fixed checkout. Original candidate key: stacks:formal-defos.tex:lemma-homotopy.\par\medskip
\noindent\textit{Editorial note.} Reference labels in the proof resolve to navigation pointers to the original author source; these pointers do not reproduce proofs of the prerequisites.\par\medskip
\noindent\textit{Allowed standard prerequisites (editorial).}\begin{itemize}
\item Prerequisites named by the author inside the counted span. The counted span is reproduced verbatim, including the author's own citations. Its prerequisites are the earlier results of the same fixed Stacks Project version that the author cites there and in the chapter's lead-in.
\end{itemize}\medskip
\ifdefined\hypertarget\hypertarget{sourcepointer-1}{}\fi
\noindent\textit{Source reference.} Original lemma 06SB: sources/formal-defos.tex, source lines 906--906. Navigation pointer used in the counted span; the span contains the pointer only and the cited result is not reproduced inside it. Registered as a non-load-bearing navigation stub in the established form for Stacks material. No bibliography citation is involved.\par\medskip
\begin{lemma}
\label{lemma-homotopy}
Being formally homotopic is an equivalence relation on
sets of morphisms in $\widehat{\mathcal{C}}_\Lambda$.
\end{lemma}

\begin{proof}
Suppose we have any $r \geq 1$ and two maps
$h_1, h_2 : R \to R[[t_1, \ldots, t_r]]$ such that
$h_1(a) \bmod (t_1, \ldots, t_r) =  h_2(a) \bmod (t_1, \ldots, t_r) = a$
for all $a \in R$ and a map $k : R[[t_1, \ldots, t_r]] \to S$.
Then we claim $k \circ h_1$ is formally homotopic to $k \circ h_2$.
The symmetric inherent in this claim will show that our notion
of formally homotopic is symmetric. Namely, the map
$$
\Psi :
R[[t_1, \ldots, t_r]]
\longrightarrow
R[[t_1, \ldots, t_r]],\quad
\sum a_I t^I \longmapsto \sum h_1(a_I)t^I
$$
is an isomorphism. Set $h(a) = \Psi^{-1}(h_2(a))$ for $a \in R$ and
$k' = k \circ \Psi$, then we see that $(r, h, k')$ is a formal
homotopy between $k \circ h_1$ and $k \circ h_2$, proving the claim

\medskip\noindent
Say we have three maps $f_1, f_2, f_3 : R \to S$ as above
and a formal homotopy $(r_1, h_1, k_1)$ between $f_1$ and $f_2$
and a formal homotopy $(r_2, h_2, k_2)$ between $f_3$ and $f_2$ (!).
After relabeling the coordinates we may assume
$h_2 : R \to R[[t_{r_1 + 1}, \ldots, t_{r_1 + r_2}]]$
and
$k_2 : R[[t_{r_1 + 1}, \ldots, t_{r_1 + r_2}]] \to S$.
By choosing a suitable isomorphism
$$
R[[t_1, \ldots, t_{r_1 + r_2}]]
\longrightarrow
R[[t_{r_1 + 1}, \ldots, t_{r_1 + r_2}]]
\widehat{\otimes}_{h_2, R, h_1}
R[[t_1, \ldots, t_{r_1}]]
$$
we may fit these maps into a commutative diagram
$$
\xymatrix{
R \ar[r]_{h_1} \ar[d]^{h_2} &
R[[t_1, \ldots, t_{r_1}]] \ar[d]^{h_2'} \\
R[[t_{r_1 + 1}, \ldots, t_{r_1 + r_2}]] \ar[r]^{h_1'} &
R[[t_1, \ldots, t_{r_1 + r_2}]]
}
$$
with $h_2'(t_i) = t_i$ for $1 \leq i \leq r_1$ and
$h_1'(t_i) = t_i$ for $r_1 + 1 \leq i \leq r_2$.
Some details omitted.
Since this diagram is a pushout in the category
$\widehat{\mathcal{C}}_\Lambda$ (see proof of
Lemma \ref{lemma-CLambdahat-pushouts})
and since $k_1 \circ h_1 = f_2 = k_2 \circ h_2$ we conclude
there exists a map
$$
k : R[[t_1, \ldots, t_{r_1 + r_2}]] \to S
$$
with $k_1 = k \circ h_2'$ and $k_2 = k \circ h_1'$.
Denote $h = h_1' \circ h_2 = h_2' \circ h_1$.
Then we have
\begin{enumerate}
\item $k(h_1'(a)) = k_2(a) = f_3(a)$, and
\item $k(h_2'(a)) = k_1(a) = f_1(a)$.
\end{enumerate}
By the claim in the first paragraph of the proof this shows that
$f_1$ and $f_3$ are formally homotopic.
\end{proof}

\end{document}
